% ===== 个人简历 · 鲁振哲 =====
% 编译: xelatex resume.tex
\documentclass[UTF8,11pt,a4paper]{ctexart}

% ----- Packages -----
\usepackage{geometry}
\geometry{left=2cm,right=2cm,top=1.8cm,bottom=1.8cm}

\usepackage{xcolor}
\definecolor{prussian}{HTML}{2E4A62}
\definecolor{oxide}{HTML}{C44536}
\definecolor{steel}{HTML}{5A7D9A}
\definecolor{ink}{HTML}{1C2838}

\usepackage{hyperref}
\hypersetup{
  colorlinks=true,
  linkcolor=prussian,
  urlcolor=oxide,
  pdftitle={鲁振哲 - 个人简历},
  pdfauthor={鲁振哲}
}

\usepackage{enumitem}
\setlist{nosep,leftmargin=18pt,itemsep=2pt}
\setlist[itemize]{label=\textcolor{oxide}{\small\textbullet}}

\usepackage{tabularx}
\usepackage{booktabs}
\usepackage{parskip}
\setlength{\parskip}{4pt}

% ----- Custom Commands -----
\newcommand{\sectionline}[1]{%
  \par\vspace{10pt}
  \noindent{\color{prussian}\large\heiti #1}
  \par\vspace{2pt}\noindent{\color{prussian}\rule{\textwidth}{1.5pt}}
  \vspace{6pt}
}

\newcommand{\cvdate}[1]{{\color{steel}\footnotesize\kaishu #1}}
\newcommand{\cvtag}[1]{%
  {\sffamily\fontsize{7.5}{9}\selectfont\bfseries%
   \textcolor{white}{\colorbox{prussian}{\rule{0pt}{8pt}\hspace{3pt}#1\hspace{3pt}}}}%
}
\newcommand{\cvmono}[1]{{\color{steel}\ttfamily\fontsize{8.2}{10}\selectfont #1}}

\begin{document}
\thispagestyle{empty}

% ===== HEADER =====
\begin{minipage}[t]{0.62\textwidth}
  {\Huge\heiti 鲁振哲}\\[2pt]
  {\large\color{steel}\kaishu Zhenzhe Lu}\\[8pt]
  \small
  \color{prussian}浙江大学管理学院 \quad
  \color{steel}信息管理与信息系统\\[2pt]
  \color{steel}本科在读 $\cdot$ 2024.09 --- 2028.06
\end{minipage}
\hfill
\begin{minipage}[t]{0.35\textwidth}
  \raggedleft\small
  \href{mailto:luzhenzhe2006@qq.com}{luzhenzhe2006@qq.com}\\[3pt]
  131 9968 3838\\[3pt]
  浙江 $\cdot$ 杭州 $\cdot$ 浙大紫金港\\[3pt]
  \href{https://gitee.com/lu-zhenzhe/personal}{\ttfamily\fontsize{8.5}{10}\selectfont gitee.com/lu-zhenzhe/personal}
\end{minipage}

\vspace{6pt}
\par\noindent{\color{prussian}\rule{\textwidth}{2.5pt}}

% ===== EDUCATION =====
\sectionline{教育背景}

\noindent
\begin{tabularx}{\textwidth}{@{}Xr@{}}
  \textbf{浙江大学（Zhejiang University）} & \cvdate{2024.09 -- 至今} \\
  \multicolumn{2}{@{}l@{}}{管理学院 \textbar\ 信息管理与信息系统 \textbar\ 本科在读} \\
  \multicolumn{2}{@{}p{\textwidth}@{}}{\footnotesize\color{steel}主修：数据结构、数据库系统、运筹学、数据分析、管理信息系统、会计学、经济学}
\end{tabularx}

\vspace{6pt}
\noindent
\begin{tabularx}{\textwidth}{@{}Xr@{}}
  \textbf{龙江县第一中学} & \cvdate{2021.09 -- 2024.06} \\
  \multicolumn{2}{@{}l@{}}{\footnotesize\color{steel}理科高中，培养了扎实的数学和逻辑思维能力}
\end{tabularx}

% ===== PROJECTS =====
\sectionline{项目经验}

% --- Project 1 ---
\noindent
\begin{tabularx}{\textwidth}{@{}Xr@{}}
  {\large\textbf{钱学森问题扩展求解：绕日返回轨道设计}} & \cvdate{2026.06} \\
  \multicolumn{2}{@{}l@{}}{%
    \cvtag{独立完成}\hspace{6pt}%
    \cvmono{Python | NumPy | Matplotlib | AstroPy | LaTeX | GA | N体模拟}%
  }
\end{tabularx}

\vspace{2pt}
\noindent\footnotesize
基于钱学森《星际航行概论》框架，开发高精度N体积分器与多重打靶优化算法，实现包含月球借力与发射窗口优化的绕日返回轨道设计。

\begin{itemize}
  \item 自研 Velocity-Verlet 辛积分器，能量漂移低至 $4.32 \times 10^{-15}$（1年仿真）
  \item 复现钱学森拼接圆锥曲线算例，所有参数偏差 $\leq 0.1\%$
  \item 遗传算法 + 并行蒙特卡洛全局优化框架，对接 JPL Horizons 真实历表残差校验
\end{itemize}

% --- Project 2 ---
\noindent
\begin{tabularx}{\textwidth}{@{}Xr@{}}
  {\large\textbf{SeekBook 二手书与知识资产流转系统}} & \cvdate{2026.05} \\
  \multicolumn{2}{@{}l@{}}{%
    \cvtag{主要开发者}\hspace{6pt}%
    \cvmono{Vue 3 | Flask | SQLAlchemy | SQLite | Qwen-VL | ECharts | Vite}%
  }
\end{tabularx}

\vspace{2pt}
\noindent\footnotesize
Vue 3 + Flask 全栈Web应用，十万级虚拟书籍数据支撑的二手书交易与知识社区平台。

\begin{itemize}
  \item 设计并实现 Blueprint 微服务工厂架构，单文件严格 $\leq$ 200行
  \item 开发 ISBN 双轨智能识别：通义千问 Qwen-VL 云端OCR + 文件名后门
  \item 构建知识星河社区模块（笔记发布、模糊搜索关联、评论互动）
  \item 实现课表 JSON 导入 $\rightarrow$ 智慧清单生成 $\rightarrow$ ECharts 知识星图可视化
  \item 线下对交交易模式：零余额系统，见面信息 SQLite 强持久化存储
\end{itemize}

% ===== RESEARCH =====
\sectionline{研究方向}

\noindent
\begin{minipage}[t]{0.32\textwidth}
  \centering\small
  \textbf{数据分析}\\[2pt]
  {\footnotesize\color{steel}统计学习、机器学习在\\商业决策中的应用}
\end{minipage}
\hfill
\begin{minipage}[t]{0.32\textwidth}
  \centering\small
  \textbf{运筹优化}\\[2pt]
  {\footnotesize\color{steel}数学规划、优化算法\\供应链与资源分配}
\end{minipage}
\hfill
\begin{minipage}[t]{0.32\textwidth}
  \centering\small
  \textbf{数学建模}\\[2pt]
  {\footnotesize\color{steel}跨学科建模挑战\\实际问题求解}
\end{minipage}

% ===== SKILLS =====
\sectionline{技能}

\noindent
\begin{tabularx}{\textwidth}{@{}p{2.8cm}X@{}}
  \textbf{编程语言} & \cvmono{Python | SQL | JavaScript | HTML/CSS} \\[4pt]
  \textbf{框架与工具} & \cvmono{Vue 3 | Flask | Vite | Git | LaTeX} \\[4pt]
  \textbf{数据分析} & \cvmono{NumPy | Pandas | Matplotlib} \\[4pt]
  \textbf{语言能力} & \cvmono{CET-4 / CET-6 | 普通话}
\end{tabularx}

% ===== AWARDS =====
\sectionline{获奖情况}

\begin{itemize}
  \item 全国大学生数学竞赛（非数学类B类）\hspace{6pt}\cvtag{浙江省一等奖}
  \item 微积分竞赛\hspace{6pt}\cvtag{浙江省三等奖}
  \item 浙江大学丹青学园团总支篮球赛\hspace{6pt}\cvtag{团体第一名}
  \item CET-4 / CET-6 英语等级测试\hspace{6pt}\cvtag{已通过}
\end{itemize}

% ===== STUDENT WORK =====
\sectionline{学生工作}

\noindent
\begin{tabularx}{\textwidth}{@{}Xr@{}}
  \textbf{管理学院学业指导中心 \textbar\ 干事} & \cvdate{2024.09 -- 至今}
\end{tabularx}

\vspace{2pt}
\noindent\footnotesize\color{steel}
参与管理学院学业指导中心的日常运营与活动组织工作，协助开展学业帮扶、经验分享等活动。

% ===== SELF-EVALUATION =====
\sectionline{自我评价}

\noindent\small
本科在读，主修信息管理与信息系统专业。关注数据分析、运筹优化与数学建模方向，具备扎实的数学基础和逻辑思维能力。熟悉 Python 科学计算与 Vue 全栈开发，具备独立完成复杂项目的能力。善于沟通协作，乐于探索新技术，致力于将数据驱动方法应用于实际管理问题。

\vspace{6pt}
\begin{center}
  \small\color{prussian}
  所有项目源码托管于\ \href{https://gitee.com/lu-zhenzhe/personal}{\textbf{\ttfamily gitee.com/lu-zhenzhe/personal}}
\end{center}

% ===== FOOTER =====
\vspace{8pt}
\par\noindent{\color{prussian}\rule{\textwidth}{1pt}}
\vspace{4pt}
\begin{center}
  {\footnotesize\color{steel}鲁振哲 \textbar\ Zhenzhe Lu \textbar\ 浙江大学管理学院 \textbar\ \today}
\end{center}

\end{document}